% 第 55 章　光与原子的合作
\chapter{光与原子的合作}

\step{反常识开场}

\begin{mainline}
拿一支最普通的红色激光笔——文具店几块钱的那种——再拿一支手电筒。两者用差不多的电池，功率都在毫瓦到瓦的量级。晚上把它们同时对准二十米外的墙：手电筒的光摊成一张直径一两米、边缘模糊的“大饼”；激光笔却只是一个指甲盖大的亮斑，亮得刺眼。再比比“颜色”：手电筒的白光用旧光盘一照，能展开成一条完整的彩虹；激光笔的红却纯得几乎没有杂质。打个比方：激光像一支全体只唱一个音、踩同一个节拍的合唱团，手电筒像菜市场里千人同时喊话。这个比方像的地方，是“整齐”与“杂乱”的对比；不像的地方在于，真人合唱再练也会走调，而激光的整齐不是靠排练，是靠一条物理定律硬拉出来的。

第二个现象可能就在你的卧室里：天花板上的夜光星星贴纸。开灯照一会儿，关灯之后它还能幽幽地亮上十几分钟。光居然能像零钱一样先存起来、再慢慢花掉。更奇怪的是“喂”它的方式：你用又亮又刺眼的红光去照它，它几乎“不吃”；换成昏暗得多的紫光一照，它立刻“吃饱”。亮光居然输给了暗光。

这两件小事背后是同一套规则：原子收光、放光的方式，比我们想象的挑剔得多，也可以被组织得整齐得多。其实你早就见过这套规则的更多演出：验钞灯一照，纸币上浮出平时看不见的荧光图案；舞厅的紫光灯下，白 T 恤亮得发蓝；加油站罩棚上的钠灯、路口的红绿灯、手机闪光灯里的白光 LED——每一种人造光源，都是原子在按自己的规矩收光、放光。本章就来讲光与原子怎样合作——从荧光、激光一路讲到光谱学和原子钟。
\end{mainline}

\step{先猜一猜}

\begin{mainline}
往下读之前，先把自己的直觉押在桌面上，别急着看答案：

\begin{itemize}[leftmargin=1.8em]
  \item 激光颜色那么纯，你猜是因为里面加了高级滤镜？还是因为发光材料天生只发一种颜色？还是另有原因？
  \item 激光几乎不散开，你猜是因为发光点排得特别整齐吗？
  \item 夜光贴纸“存光”，你猜它像一面反射得很慢的镜子，还是像一块能充电的电池？
  \item 红光充不进贴纸、紫光充得进，你猜关键差别是光的“总量”，还是光的“单价”？
\end{itemize}

把你的答案写在纸上。本章结尾可以回来对账。提醒一句：在这类问题上下注，物理学家的直觉也经常输——二十世纪初最聪明的一批人，就曾集体押错。他们原以为只要把电磁学算得再细一点，原子的账本总会平衡；结果逼着他们改账本的，正是你现在桌上那支激光笔的祖先们。
\end{mainline}

\step{动手实验}

\begin{experiment}[光盘分光与荧光水]
准备：一张旧 CD 或 DVD、一支红色激光笔、一支黄色荧光笔、一个透明杯子、一张夜光贴纸（文具店或玩具店有售）；如果家里有验钞紫光灯或紫色 LED 更好。安全提醒：激光绝不对着人眼，也小心镜子和瓷砖的反射。

第一步，造一台粗糙的光谱仪。关掉大灯，让一盏白炽灯（或卤素灯）的光照到 CD 的反光面上，慢慢转动光盘，用眼睛寻找反射出的彩虹；再换成 LED 台灯或节能灯重复一次。你会发现：白炽灯展开成一条连续不断的彩虹，LED 灯和节能灯却只出现几条分立的亮带，带与带之间是黑的。同一张光盘，两种灯，给出两种完全不同的“光的身份证”。

第二步，让光在水杯里“换颜色”。把黄色荧光笔的笔芯泡进半杯水里搅匀，得到一杯淡黄绿色的水。在黑暗处让红色激光从侧面水平射入杯中：水里会出现一条明亮的绿色光路。红光进去，绿光出来——光穿过这杯水时，颜色被换掉了。

第三步，比较“充电”效果。把夜光贴纸分两个区域：一个区域用红色激光或红色 LED 近距离照三十秒，另一个区域用紫光（验钞灯、紫色 LED，或者用窗外阳光）照三十秒，然后关灯比较。你会发现紫光那一角明显更亮；红光照得再久，也几乎充不进去。

请记下三条观察：光谱是连续的还是分立的；荧光的颜色与入射光相比谁更偏蓝；哪种光能给贴纸充电。这三条记录，就是本章全部理论的实验地基。
\end{experiment}

\step{旧模型遇到麻烦}

\begin{mainline}
在第 51 章之前，我们手里的旧模型长这样：光是连续流淌的电磁波，能量想多细就可以分多细；原子是一个微小的太阳系，电子绕核转圈。拿这套账本去记刚才的实验，立刻有四处对不上账。

第一处，账本本身已经破产。第 51 章算过：按照电磁学规律，绕核转圈的电子应当不断辐射电磁波、损失能量，在远小于一秒的时间内螺旋坠进原子核。可原子稳定得很。所以“微小太阳系”这个模型早就不能用了，电子没有那样一条确定的路线。

第二处，连续彩虹对不上分立亮带。灼热的固体（灯丝）发出连续光谱；而稀薄气体发光——霓虹灯、节能灯里的汞蒸气、烟花里的锶和铜——只发出若干条颜色严格固定的细线。每种元素一套，互不重复，像指纹一样。你在家就能验证一条：往燃气灶的蓝火苗上撒一小撮食盐，火苗立刻变成明亮的黄色——那是钠原子只肯唱的那个音；老城区昏黄的钠路灯，唱的也是同一个音。连续流淌、想要多细有多细的光，解释不了“只唱几个固定音”。

第三处，荧光变色。红光进、绿光出。如果光只是一股笼统的“能量流”，颜色不该这样单向地改变。你从来没有见过绿光进去、蓝光出来的贴纸——方向反过来的事，自然界不做。

第四处，充电时的“单价歧视”。刺眼的红光充不进夜光贴纸，昏暗的紫光却可以。如果能量只看总量，这是荒谬的；除非原子收能量时根本不看总量，只看每一份的“面额”——面额不对，再多零钱也买不了这张票。

再加上第 42 章里黑体辐射和光电效应已经逼出来的结论，出路只剩一条：原子的能量是一级一级分立的，光与原子交换能量只能整份整份地换，每份的面额由光的颜色决定。这就是本章要发明的新概念。
\end{mainline}

\step{发明新概念}

\begin{mainline}
新概念一：能级。原子的内部能量像楼梯，不像斜坡。它像楼梯的地方在于：只能站在第几级上，能量只能取分立的数值，从一级到另一级必须整个跳过去。它不像楼梯的地方在于：真楼梯上你可以扶着墙停在两级中间，原子却不存在“两级之间”的状态，连一瞬间都不停留；也千万别把它想成行星轨道——第 52 章讲过，电子在原子里的状态没有确定的路线，“第几级”只标记能量，不标记位置。

新概念二：吸收与自发辐射。低能级的原子碰到一个面额正好等于两级能量之差的光子，可以整个吞掉它，跳上去——这叫吸收。少一分不行，多一分也不要：不收零钱，也不找零。高处的原子待不稳，会随机地、不可预测地在某一刻跳下来，把能量差打包成一个光子吐出去——这叫自发辐射。霓虹灯的橙红色、烟花里锶的红和铜的蓝绿，就是各种原子自发辐射时的“固定找零面额”。单个原子哪一刻跳下，原则上无法预言（这与第 45 章的概率是同一副面孔）；但亿万个原子的统计规律非常稳定，稳定到可以拿来做钟，后面会讲。

新概念三：受激辐射。1916 到 1917 年，爱因斯坦为了解释热辐射的平衡，提出还存在第三种过程：一个面额恰好合适的光子路过一个激发态原子时，除了被低能级原子吸收的可能之外，还会“邀请”这个激发态原子立刻跳下，放出第二个光子。关键在于：新光子与路过的光子在频率、方向、相位上完全一致，像一份完美克隆。它像雪崩：一个雪球滚过，触发坡上的积雪跟着滚下，越滚越大。它不像雪崩的地方在于：雪球之间是机械推挤，而受激辐射并不是光子“撞”原子、撞出另一个光子；倒更像合唱指挥给出一个起拍，全体在同一拍上开口——而且克隆出的光子连一丝走调都没有，比任何真人合唱都整齐。这里要分清账本的两栏：实验事实是激光确实存在、光谱线确实分立；模型解释是“能级跃迁加光子克隆”这套记账方式，它是目前最简洁、从未失手的一本账。
\end{mainline}

\begin{figure}[htbp]
\centering
\begin{tikzpicture}[>=Stealth,scale=0.95]
% ---- 吸收 ----
\draw[thick] (0,0)--(1.8,0); \draw[thick] (0,1.5)--(1.8,1.5);
\fill (0.9,0) circle (1.6pt);
\draw[decorate,decoration={snake,amplitude=1.2pt,segment length=5pt},->] (-0.7,-0.6)--(0.75,1.2);
\node at (0.9,-1.0) {吸收};
% ---- 自发辐射 ----
\draw[thick] (3.4,0)--(5.2,0); \draw[thick] (3.4,1.5)--(5.2,1.5);
\fill (4.3,1.5) circle (1.6pt);
\draw[->,dashed] (4.3,1.35)--(4.3,0.15);
\draw[decorate,decoration={snake,amplitude=1.2pt,segment length=5pt},->] (4.5,0.75)--(6.0,-0.2);
\node at (4.3,-1.0) {自发辐射};
% ---- 受激辐射 ----
\draw[thick] (7.0,0)--(8.8,0); \draw[thick] (7.0,1.5)--(8.8,1.5);
\fill (7.9,1.5) circle (1.6pt);
\draw[decorate,decoration={snake,amplitude=1.2pt,segment length=5pt},->] (6.2,0.9)--(7.6,0.9);
\draw[decorate,decoration={snake,amplitude=1.2pt,segment length=5pt},->] (8.1,1.05)--(9.7,1.05);
\draw[decorate,decoration={snake,amplitude=1.2pt,segment length=5pt},->] (8.1,0.7)--(9.7,0.7);
\node at (7.9,-1.0) {受激辐射};
\end{tikzpicture}
\caption{光与原子的三种基本交易：黑点代表原子所处的能级，波浪箭头是光子。受激辐射中，一个光子路过，两个一模一样的光子离开。}
\end{figure}

\begin{mainline}
新概念四：粒子数反转。麻烦在于，同一种“邀请”机制对低能级原子表现为吸收，对高能级原子表现为受激辐射，两者完全对称。平时热平衡的状态下，低能级总是“人多”，光穿过介质时被吸收的比被放大的多，整体是衰减的。要让放大占上风，就必须把更多原子搬到高能级，让“楼上比楼下人多”——这种反常状态叫粒子数反转。手段叫泵浦：用电流，或者用另一束光，像用水泵把水一车车抽到山顶水库。它像水库蓄水的地方在于：先把能量存到高处，用时再放；它不像永动机的地方在于：每一份能量都来自电源，账上一分不多。具体技巧是借道第三级台阶：泵浦把原子快速送上一个短命的高能级，它立刻跌到一个能待很久的“中转站”（亚稳能级）里堆积起来，从而对下面那一级形成反转。

新概念五：荧光与磷光。荧光材料吸收一个紫外光子跳到高处，先在分子内部“换乘”几站，把一小部分能量变成振动发热，再从较低的一级跳下，放出可见光子。出来的光子面额变小了，颜色向红端挪：所以紫光进、蓝光出，红光进、绿光出。反过来蓝光进、紫光出几乎不发生——那等于让小面额合成大面额，差额没有出处。生活里的荧光比比皆是：验钞灯下纸币浮出的防伪图案、白 T 恤里让白更白的荧光增白剂、马路施工服上晃眼的荧光条，都是这套“换乘找零”机制。夜光贴纸（掺稀土的铝酸锶）多一个“停车场”：电子跳下之前先落入一个按规则很难直接下来的亚稳级，只能慢慢地漏，于是发光被拖成几分钟——这就是磷光。停车场的存在并不违反任何规则，只是某些“下楼的楼梯”被量子规则设成了限行车道（选择定则，见第 52 章）。

新概念六：发光二极管。第 54 章说过，固体里亿万个原子挤在一起，原先分立的能级摊成了能带。LED 里面，电子从导带跌进价带，一步放出一个光子，颜色由带隙的宽度决定——所以红色 LED 和蓝色 LED 用的是不同的半导体材料。家里照明用的白光 LED，其实是蓝光 LED 加一层荧光粉：蓝光泵浦，荧光粉把一部分蓝光换成黄光，蓝加黄在眼睛里配成白。你家的 LED 灯泡里，就住着本章讲的荧光。
\end{mainline}

\begin{figure}[htbp]
\centering
\begin{tikzpicture}
% 热平衡
\draw[thick] (0,0)--(2.0,0); \draw[thick] (0,1.3)--(2.0,1.3);
\fill (0.4,0) circle (1.6pt); \fill (0.8,0) circle (1.6pt); \fill (1.2,0) circle (1.6pt); \fill (1.6,0) circle (1.6pt);
\fill (1.0,1.3) circle (1.6pt);
\node at (1.0,-0.6) {热平衡：楼下人多};
% 反转
\draw[thick] (4.0,0)--(6.0,0); \draw[thick] (4.0,1.3)--(6.0,1.3);
\fill (5.0,0) circle (1.6pt);
\fill (4.4,1.3) circle (1.6pt); \fill (4.8,1.3) circle (1.6pt); \fill (5.2,1.3) circle (1.6pt); \fill (5.6,1.3) circle (1.6pt);
\node at (5.0,-0.6) {粒子数反转：楼上人多};
\end{tikzpicture}
\caption{每个黑点代表一个原子。左图是热平衡的常态，吸收占上风；右图是泵浦制造出的反转，受激辐射可以雪崩式放大。}
\end{figure}

\begin{mainline}
把这本新账翻过来用，就是光谱学——物理学里最会“读心”的技术。每种元素的能级表独一无二，所以发光时给出的谱线就是它的指纹；反过来，让一束连续光谱的光（比如恒星表面的热辐射）穿过较冷的气体，气体里的原子只吞掉自己面额对得上的光子，光谱上就出现一条条暗线——这叫吸收线。亮线和暗线，是同一本能级账的收、支两栏。

光谱学的战利品清单长得惊人。1868 年日全食，天文学家在太阳色球的光谱里发现一条当时地球上没人认识的黄色亮线，断定那属于一种未知元素，用希腊语“太阳”命名，就是氦；二十多年后，氦才在地球上的矿物里被找到——人类先在太阳上、后在地球上发现了一种元素。今天，炼钢厂的直读光谱仪把钢水溅出的火花一照，几秒钟就报出碳、锰、铬的含量，师傅据此决定要不要加料；刑侦实验室用同样的办法比对油漆和玻璃的微量元素；天文台则靠恒星光谱读出它的成分、表面温度，再靠整组谱线的红移或蓝移读出它靠近或远离我们的速度。

激光自己也是工程界的主力军：超市收银台的条形码扫描、入户宽带里的光纤通信（一根头发丝粗细的玻璃丝里跑着激光）、工厂里的激光切割与焊接、近视手术用的准分子激光——那种紫外激光每个光子的面额大到可以直接打断角膜分子键，组织还来不及发热就被“气化”掉一层，所以切口几乎不烫伤周围。这些应用的共同点只有一个：把光的面额、方向或节拍用到极致。
\end{mainline}

\step{一句话模型}

\begin{oneline}
原子只按固定面额收发能量包裹，面额由颜色决定：荧光是“收大面额、找零成小面额慢慢花”，激光是“先用泵浦制造粒子数反转，再让受激辐射把亿万个包裹的面额、方向和节拍统一成同一个”。
\end{oneline}

\step{数学翻译器}

\begin{mainline}
面额公式只有一条：
\[
\Delta E=h\nu=\frac{hc}{\lambda}
\]
左边回答的问题是“这级台阶有多高”；右边每一项都在改变台阶：\(\nu\) 是光的频率，\(\lambda\) 是波长，\(h\) 是普朗克常量，\(h\approx 6.63\times10^{-34}\,\mathrm{J\cdot s}\)，\(c\approx 3.0\times10^{8}\,\mathrm{m/s}\)。拿红色激光笔练手：\(\lambda\approx 650\,\mathrm{nm}\)，于是 \(\nu=c/\lambda\approx 4.6\times10^{14}\,\mathrm{Hz}\)，每个光子的面额 \(E\approx 3.1\times10^{-19}\,\mathrm{J}\)。这个数字太小，工程上常用电子伏：\(1\,\mathrm{eV}\approx 1.6\times10^{-19}\,\mathrm{J}\)，所以红光光子约 \(1.9\,\mathrm{eV}\)，紫光（\(\lambda\approx 400\,\mathrm{nm}\)）光子约 \(3.1\,\mathrm{eV}\)。

立刻做特殊情形体检。频率趋于零（波长极长的无线电波），面额趋于零——所以电台天线发射的能量虽大，单个“包裹”却小到几乎不扰动原子，无线电波照在身上不会“充电”。波长减半，面额翻倍——这就是为什么紫外线能晒伤皮肤、能给贴纸充电，可见光不能。方向反转也对得上：吸收一个 \(1.9\,\mathrm{eV}\) 光子跳上去的原子，沿同一对台阶跳下来时只能吐回 \(1.9\,\mathrm{eV}\)，一进一出账永远平。

真实数据估算：一支 \(5\,\mathrm{mW}\) 的红激光笔，每秒发出多少个光子？
\[
N=\frac{P}{E}=\frac{0.005\,\mathrm{J/s}}{3.1\times10^{-19}\,\mathrm{J}}\approx 1.6\times10^{16}\ \text{个/秒}
\]
体检：\(P\) 趋于零则 \(N\) 趋于零，合理；同样的 \(5\,\mathrm{mW}\) 换成紫外激光，面额更大，光子数更少，也合理。每秒一万六千亿个光子听起来吓人，但每个面额太小，所以激光笔烧不穿纸。可它又足以让远处的眼睛一个一个地收：人眼在最暗适应时，几百个光子落进瞳孔就足以感知。

再看荧光的颜色账。贴纸吸收一个 \(400\,\mathrm{nm}\) 紫光光子（\(3.1\,\mathrm{eV}\)），内部换乘发热约 \(0.4\,\mathrm{eV}\)，放出约 \(2.7\,\mathrm{eV}\) 的蓝光。出射光子的能量总是小于等于入射光子，这条规律叫斯托克斯规则。它的零情形是：换乘损失为零时，出射等于入射，那就不再叫荧光，而是普通的散射。
\end{mainline}

\begin{mathbridge}
为什么热平衡下永远造不出粒子数反转？温度为 \(T\) 的平衡体系里，高能级与低能级上的原子数目之比是
\[
\frac{N_2}{N_1}=e^{-\Delta E/(kT)}
\]
这是玻尔兹曼分布（第 33 章）的直接推论，\(k\approx 1.38\times10^{-23}\,\mathrm{J/K}\)。体检三个极限：\(T\to 0\)，指数趋于负无穷，比值趋于零，全部原子都在楼下，合理；\(T\to\infty\)，指数趋于零，比值趋于一——最多相等，永远到不了“楼上人多”，所以单纯加热这条路造不出反转；室温下 \(kT\approx 1/40\,\mathrm{eV}\)，对 \(\Delta E\approx 2\,\mathrm{eV}\) 的可见光台阶，比值约为 \(e^{-80}\approx 10^{-35}\)。这就是室温物体不会自己发可见光的原因；灯丝要烧到两三千度，才从热辐射里挤出一点可见光。结论：反转必须靠“挑选手法”的泵浦——电流、另一束光、化学反应——绕开热平衡的限制。

三能级激光的速率图像（只读方程，不求解）：泵浦以速率 \(R\) 把原子从 1 级送上 3 级；3 级极短命，原子几乎立刻无辐射弛豫到亚稳的 2 级；2 级以平均寿命 \(\tau\) 慢慢漏回 1 级。于是
\[
\frac{dN_2}{dt}=R-\frac{N_2}{\tau}-(\text{受激辐射项})
\]
把 \(R\) 开大，\(N_2\) 就堆积起来，对 1 级形成反转；而受激辐射项随腔内光强增大而增大——光越强，邀请越多，产生的光更强。这正是“放大”的自我强化回路，也是激光器起振的数学骨架。
\end{mathbridge}

\begin{challenge}
爱因斯坦最初的表述用三个系数：自发辐射率 \(A_{21}\)，受激辐射率 \(B_{21}\rho(\nu)\)，吸收率 \(B_{12}\rho(\nu)\)，其中 \(\rho(\nu)\) 是辐射场在频率 \(\nu\) 附近的能量密度。要求这套速率与普朗克黑体辐射谱在热平衡下相容，会强制推出 \(B_{12}=B_{21}\)，以及 \(A_{21}/B_{21}\) 正比于 \(\nu^{3}\)。换句话说，受激辐射的存在不是一个额外的假设，而是“原子加热辐射平衡”这个体系逼出来的逻辑必然——1917 年的爱因斯坦是在激光器诞生四十多年前，从账本平衡里先看见了它。

再看谐振腔。光在增益介质中单程走过长度 \(L\)，强度被放大 \(e^{gL}\) 倍，\(g\) 叫增益系数；两面镜子的反射率小于一，每往返一次还要漏掉一部分。起振条件写成“单程增益恰好抵消单程损耗”。阈值以下，腔内只有被介质稍稍过滤的微弱荧光；阈值以上，满足驻波条件的某一个模式雪崩式胜出，其余模式被“抢”光了反转粒子而饿死。从这个角度看，一台激光器是一个有相变味道的非平衡系统：泵浦功率就是那个被缓缓拧过临界点的旋钮。
\end{challenge}

\begin{mainline}
同一个公式还解释了一项你最熟悉的工程：白光 LED 为什么省电。白炽灯靠把灯丝烧到近三千度，靠热辐射“挤”出可见光，大部分能量变成红外线白白发热；LED 里电子直接跨带隙跳下，一步一个光子，能量几乎全用在发光上。面额公式在这里换了副面孔：不是“台阶有多高”，而是“这份电能变成了多少个包裹、每个包裹值多少”。

最后是精确的极限：原子钟。现行国际单位制里，一秒的定义是铯-133 原子基态两个超精细能级之间跃迁所对应辐射的 9192631770 个周期。为什么拿原子当钟，就能准到几千万年才差一秒？三个原因。第一，全世界每一个铯原子的这条谱线都一模一样——不需要把一件“标准原器”锁在保险柜里，标准本身遍地都是。第二，跃迁面额由量子规则钉死，不像钟摆会热胀冷缩、石英晶体会老化。第三，频率是人类能测得最准的物理量：数周期比量长度、称质量都容易做得精细。

原子钟最出名的用户是卫星导航。定位的本质是测量“信号从卫星飞到手机用了多久”，乘以光速得到距离。光速约每秒三十万千米，所以时钟每差 1 纳秒（十亿分之一秒），距离就差约 30 厘米。卫星上的原子钟每天因为狭义和广义相对论的效应净快约 38 微秒，如果不按第 39 章和第 41 章的公式修正，导航误差每天会累积约十千米。也就是说，你手机里的每一次打车定位，都是原子钟和相对论在后台联合值班。
\end{mainline}

\step{模型的边界}

\begin{mainline}
这套模型非常好用，但每一句话都有边界，值得逐条钉在墙上。

“激光是单色光”——不是数学意义上的单一频率。任何激光都有有限的线宽，腔内原子的热运动、腔长的微小抖动都会给频率抹上一点宽度。一台稳频激光器可以做到极窄，但“极窄”不等于“零”。

“激光绝对平行”——不成立。衍射给任何光束设了下限：从直径 \(D\) 的孔径里发出的光，发散角至少是 \(\lambda/D\) 的量级。激光笔出口只有一两毫米，红光在几十米外摊成硬币大小，正是这条下限在工作，不是激光器“不够好”。

“受激辐射克隆光子”——克隆的是频率、方向、相位这些“模式属性”，不是把某个光子携带的信息复制一份再超光速送走。这里没有任何超光速传信的通道，第 48 章谈纠缠时守的是同一条规矩。

“两个能级加一束强光，就能做出激光”——不成立。同一频率的光既是泵浦又诱发受激辐射，吸收和放大对称地抵消，最好也只能做到上下能级人数相等，永远到不了反转。真实激光器至少借道第三能级，或者用电流、化学反应把激发态直接“造”出来——半导体激光器靠的就是注入电子和空穴。

“荧光是物质自己发光”——不是。荧光材料的能量全部来自先吸收的光（或者电流、化学能），出射光的总能量必然小于等于吸收的能量。白 T 恤在舞厅紫光灯下“亮得过分”，是因为环境暗，外加不可见的紫外被换成了可见蓝光，不是它凭空增产了能量。

“光到底是波还是粒子”——本章从头到尾没有让光在波和粒子之间来回变身。描述它怎样传播，用波动语言，干涉、衍射全都照常；描述它怎样与原子交换能量，用一份一份的“包裹”语言。两套账记的是同一个对象，就像同一家店既按营业额记账、又按单笔小票记账，不是店在两种形态之间切换。

还要说清量子理论与经典光学的关系：经典光学并没有被推翻。透镜成像、反射折射、干涉、衍射、偏振的每一条定律，在强光下照旧精确，因为亿万光子的平均行为正是电磁波。量子理论补上的，是交换能量的“那一瞬”和极微弱的光：把光调弱到平均每秒只有几个光子进入仪器，屏幕上仍然是一个点一个点地砸出来，可点攒得多了，干涉条纹照样浮现（第 43 章详细讲过这幅图景）。经典光学是这本量子账的“总额页”，量子理论才是它的“逐笔流水”。

还有两个典型误用。其一，以为激光器“产生”能量：能量全部来自泵浦，电光转换效率往往不到一半，剩下的变成热——大功率激光器有相当一部分体积是散热器。其二，以为白光 LED 发出“白色光子”：光子没有白色这个品种，白是多种颜色的光子混在一起、在眼睛里产生的总体感觉。
\end{mainline}

\begin{figure}[htbp]
\centering
\begin{tikzpicture}[>=Stealth]
\draw[fill=orange!15] (0,0) rectangle (5,1);
\node at (2.5,0.5) {增益介质（粒子数反转）};
\draw[line width=2.2pt] (-0.12,-0.45)--(-0.12,1.45);
\node[below] at (-0.12,-0.45) {全反镜};
\draw[line width=2.2pt] (5.12,-0.45)--(5.12,1.45);
\node[below] at (5.12,-0.45) {部分透射镜};
\draw[<->,red!70!black] (0.15,0.5)--(4.85,0.5);
\draw[->,red!70!black,line width=1.6pt] (5.12,0.5)--(6.9,0.5) node[right]{激光输出};
\draw[->,blue!70!black,line width=1.2pt] (2.5,2.1)--(2.5,1.08);
\node[above] at (2.5,2.1) {泵浦（电流或另一束光）};
\end{tikzpicture}
\caption{激光器的基本结构：泵浦制造反转，两面镜子构成谐振腔，只留下沿轴线往返的模式被反复放大，一部分从右镜透出成为激光。}
\end{figure}

\step{回头解释开场现象}

\begin{mainline}
现在回到那两支笔。手电筒的灯丝是热辐射：亿万个原子在热碰撞中各自随机地跳上跳下，方向乱、颜色杂、节拍全无，能量摊在所有方向和整段彩虹上。总功率不算小，分到每个方向、每个颜色，就所剩无几了。激光笔里是一台微型半导体激光器：电流直接把电子和空穴注入出粒子数反转，晶体两端平行的解理面就是两面镜子，构成谐振腔。光在腔内来回穿行，凡是不沿轴线、不满足驻波条件的光，都漏出腔外死掉了；活下来的那一个模式被受激辐射反复克隆放大。方向、颜色、节拍，三件事被同一套机制钉死。所以同样是几毫瓦，激光把全部“预算”集中到了极小的锥角和极窄的频宽里——不是能量更多，是组织得更好。可以粗算一笔账：一支 \(5\,\mathrm{mW}\) 激光笔的光束散角约千分之一弧度，一支 \(5\,\mathrm{W}\) 手电筒的光摊进大半个球面。功率差一千倍，立体角却差上百万倍，摊下来单位方向上的亮度，激光笔反而赢出手电筒几个数量级——这就是为什么二十米外那个小红点仍然刺眼。

夜光贴纸也结案了。紫光光子的面额够，把电子送上“高架”；高架与地面之间隔着亚稳“停车场”，电子排队慢慢下楼，每分钟放出一批绿光，于是关灯后还能亮很久。红光光子的面额不足，连高架的入口都到不了，再刺眼也没用——原子不收零钱，这一条规矩，多大的功率也救不了。

现在可以对账了：激光颜色纯，不是滤镜，不是材料天生，而是“能级定面额、谐振腔选模式、受激辐射做克隆”三件套的结果；光束不散，是腔选方向加衍射下限；贴纸存光，是亚稳停车场，不是慢反射；红光输给紫光，是单价问题，不是总量问题。你押对了几条？
\end{mainline}

\step{脑洞实验与挑战题}

\begin{thought}[月光下的几个光子]
1969 年起，阿波罗计划的宇航员在月面安放了几组角反射器。今天，地球上的天文台仍在向它们发射激光脉冲：一个脉冲含有约 \(10^{17}\) 个光子，飞行约 2.5 秒往返之后，能回到望远镜里的，平均每次只有几个光子——极端吧？单程之后，光斑已在月面上摊成几公里宽，绝大多数光子根本碰不到反射器。

可就是这几个“漏网之鱼”，让地月距离的测量精确到毫米级，并发现月球正以每年约 3.8 厘米的速度远离地球。想一想：发出去 \(10^{17}\) 个、只收回几个，这个测量为什么还成立？提示：测量不需要收回全部光子，只需要收回的光子忠实携带“飞了多久”的信息；而每一次“收回”，都是望远镜探测器里某个原子吸收一个光子、给出一次点击。光与原子的合作，一直进行到探测的最后一刻。

再追问一层：既然每个脉冲只回来几个光子，为什么不干脆把激光调亮一万倍？答案是，回来几个就已经够了——只要计数统计足够干净，几个光子给出的到达时刻依然精确。这倒过来提醒我们：真正限制测量的往往不是“信号太弱”，而是“噪声里还能不能认出信号”。把激光的节拍做纯、把探测器原子的响应做快，都是在给这条账降噪。
\end{thought}

\begin{mainline}
本章也是通往后面几章的桥。把同样的“能级加面额”规律搬进原子核，台阶的高度从几个电子伏变成几十万、几百万电子伏，包裹改名叫 \(\gamma\) 射线——那是第 56 章的故事。而受激辐射的光子“全同到不可分辨”这一条，其实就是第 53 章里玻色子性格的一块广告牌；把它追到底，会逼出现代的图景：光子是电磁场的激发，激光是同一个场模式里挤满了光子。这个图景不是书斋里的摆设：引力波探测器用一束高功率、频率稳到极致的激光，在几公里长的干涉仪两臂之间来回反射，量出比质子直径还小上千倍的臂长变化——正是“光子全同、节拍统一”这条性质，让人类听见了黑洞相撞的涟漪。光与原子的合作，到这里才刚刚开场。
\end{mainline}

\begin{puzzles}
\item 为什么商店里没有“白光激光笔”？提示：先回答激光的“纯”是被哪两件东西钉死的；再想想一个谐振腔、一套粒子数反转，能不能同时纵容所有颜色。
\item 荧光棒掰一掰就发光，萤火虫在夜里也会发光，它们并没有被灯照过。这违反本章“荧光必须先吸收光”的说法吗？提示：跟踪能量的来源——泵浦不一定非得是光；想想“先把能量存上高架”可以有哪些付款方式。
\item 验钞灯下，含荧光增白剂的白衬衫发出明亮的蓝紫光，关灯的一瞬间立刻熄灭；夜光贴纸却能亮上十分钟。两者都是“先吸收后放出”，差别出在哪一级台阶上？提示：问中间有没有“停车场”，以及停车场的出口车道有多宽。
\item 天文台收到某个遥远星系的光，发现每条谱线的位置都比实验室里氢的谱线整体向红端挪了一点，但谱线之间的相对间隔与氢的指纹完全一致。为什么天文学家敢断定“那就是氢”？这个整体红移又可能在报告什么？提示：认元素看相对位置；整体平移是另一种物理——想想第 38 章的多普勒效应和宇宙的膨胀。
\end{puzzles}
